% !TeX root = ../navier-stokes-manuscript.tex
% ============================================================
% 2.3 The Pressure-Complete Mixed Jet
% ============================================================

\subsection{The Pressure-Complete Mixed Jet}\label{sub:ThePressureCompleteMixedJet}

At temporal order \(n\), the equation keeps four pieces together:
\(V_{n+1}\), \(\nu\Delta V_n\), every temporal split of transport, and the
normalized pressure row \(Q_n\).  A single spatial derivative already shows
how this row widens:
\[
  \partial_\ell\bigl((V_a\cdot\nabla)V_{n-a}\bigr)
  =
  ((\partial_\ell V_a)\cdot\nabla)V_{n-a}
  +(V_a\cdot\nabla)\partial_\ell V_{n-a}.
\]
At multi-index order \(\alpha\), every derivative must be distributed between
the transporting and transported factors.  Denote the resulting coordinates by
\[
  J_{n,\alpha}
  :=
  \partial_t^n\partial_x^\alpha u,
  \qquad
  \Pi_{n,\alpha}
  :=
  \partial_t^n\partial_x^\alpha p,
\]
and use
\[
  \binom{\alpha}{\beta}
  :=
  \prod_{k=1}^{3}\binom{\alpha_k}{\beta_k}.
\]
The two Leibniz rules then give the complete transport split
\[
  \partial_t^n\partial_x^\alpha\bigl((u\cdot\nabla)u\bigr)
  =
  \sum_{a=0}^{n}
  \sum_{\beta\le\alpha}
  \binom{n}{a}
  \binom{\alpha}{\beta}
  (J_{a,\beta}\cdot\nabla)
  J_{n-a,\alpha-\beta}.
\]
Substitution keeps that split beside viscosity and pressure:
\begin{align*}
  J_{n+1,\alpha}
  &+
  \sum_{a=0}^{n}
  \sum_{\beta\le\alpha}
  \binom{n}{a}
  \binom{\alpha}{\beta}
  (J_{a,\beta}\cdot\nabla)
  J_{n-a,\alpha-\beta}
  \\
  &+\nabla\Pi_{n,\alpha}
  =
  \nu\Delta J_{n,\alpha},
  \qquad
  \nabla\cdot J_{n,\alpha}=0.
\end{align*}
Taking divergence supplies the corresponding pressure row:
\[
  -\Delta\Pi_{n,\alpha}
  =
  \partial_i\partial_j
  \sum_{a=0}^{n}
  \sum_{\beta\le\alpha}
  \binom{n}{a}
  \binom{\alpha}{\beta}
  (J_{a,\beta})_i
  (J_{n-a,\alpha-\beta})_j.
\]

Every displayed coordinate is a derivative of the original pressure-complete
pair.  A finite rectangle will retain only finitely many of them, but it never
changes the underlying solution.

The viscous coupling links the next temporal row to two additional spatial
derivatives of the present row.  The critical-height energy calculation isolates
that ascent while preserving the transport and pressure terms that travel with
it.
